% =====================================================================
\section{The 600-cell response substrate}\label{sec:substrate}
% =====================================================================

This section constructs the substrate. \S\ref{ssec:vertices}
gives the vertex set. \S\ref{ssec:graph} gives the graph and its
computed Laplacian spectrum. \S\ref{ssec:cphi} gives the response
operator $\Cph$ and the $\Ph^{-2}$ stability clamp.
\S\ref{ssec:shells} gives the 9-shell decomposition used for source
projection. \S\ref{ssec:cascade} sketches the seven-rung cascade
descent referenced by the recurrent layer in~\S\ref{sec:chain}. None
of these objects depend on neural data.

\subsection{Vertex construction}\label{ssec:vertices}

The 600-cell $\Rsixhundred$ has $120$ vertices in
$\mathbb{R}^{4}$~\citep{CoxeterRegularPolytopes,Weisstein600Cell}.
With $\Ph = (1+\sqrt 5)/2$ the canonical vertex set is
\begin{itemize}\itemsep=1pt
\item $8$ vertices: all permutations of $(\pm 1, 0, 0, 0)$;
\item $16$ vertices: all sign combinations of
  $(\pm 1, \pm 1, \pm 1, \pm 1)/2$;
\item $96$ vertices: all even permutations of
  $(\pm \Ph, \pm 1, \pm 1/\Ph, 0)/2$.
\end{itemize}
All $120$ vertices lie on the unit $3$-sphere $S^{3}$. The H$_4$
Coxeter group acts transitively on the vertex set; in particular,
every vertex has identical local structure. Implementation:
\texttt{kernel/vfd\_closure\_kernel.py:build\_600cell\_vertices}.

\subsection{Graph and Laplacian spectrum}\label{ssec:graph}

The substrate graph $G_{V_{600}} = (V, E)$ is built by connecting each
vertex to its nearest neighbours under the Euclidean metric on $S^{3}$.
H$_4$ acts transitively on the vertex set, forcing uniformity of the
local structure. We compute the resulting Laplacian spectrum from the
constructed graph; multiplicities match the expected H$_4$ block
sizes. The 600-cell construction itself is
standard~\citep{CoxeterRegularPolytopes,Weisstein600Cell}.
\paragraph{Graph facts (forced by the construction).}
The graph $G_{V_{600}}$ has $|V|=120$ vertices, $|E|=720$ edges, and
every vertex has degree exactly $12$ (H$_4$ transitivity acts on the
vertex set; the short-edge nearest-neighbour graph inherits this
uniformity). These facts are standard~\citep{CoxeterRegularPolytopes,Weisstein600Cell}
and reproduced numerically by
\texttt{kernel/vfd\_closure\_kernel.py:build\_600cell\_vertices}.

\paragraph{Laplacian spectrum (computed numerically).}
The unweighted graph Laplacian $\Lop = D - A$ has nine distinct
eigenvalues with multiplicities summing to $120$:
\[
\sigma(\Lop) \;=\;
\bigl\{0^{1},\;
       (12\!-\!6\Ph)^{4},\;
       (12\!-\!4\Ph)^{9},\;
       9^{16},\;
       12^{25},\;
       14^{36},\;
       (4\Ph + 8)^{9},\;
       15^{16},\;
       (6\Ph + 6)^{4}\bigr\},
\]
i.e.\ approximately $\{0, 2.292, 5.528, 9, 12, 14, 14.472, 15, 15.708\}$
with multiplicities $\{1, 4, 9, 16, 25, 36, 9, 16, 4\}$. We computed
this directly from the constructed Laplacian
(\texttt{kernel/vfd\_closure\_kernel.py:compute\_graph\_laplacian});
the spectrum is reproducible at machine precision and the
multiplicities match the expected H$_4$ block sizes. We do not derive
the closed-form entries here; the values in $\mathbb{Z}[\Ph]$ are
reported as observed.

\paragraph{H$_4$ irrep block decomposition.}
The eigenspaces partition into H$_4$-proper and $\sigma$-twin Coxeter
exponent classes. For H$_4$ proper the exponents are $\{1, 11, 19, 29\}$;
for the Galois twin $\sigma(\mathrm{H}_4)$ under the $\sigma$-automorphism
of $\mathbb{Z}[\Ph]$ the exponents become $\{7, 13, 17, 23\}$. The
$\sigma$-orbit projector basis
(\texttt{kernel/sigma\_orbit\_basis.py:\_self\_test}) realises the block
decomposition with cross-block norm $<10^{-15}$, providing a
machine-precise structural index used by the recurrent layer in
\S\ref{sec:chain} (the $K_{7}$-class projector is the default
phenomenal-binding profile).

\subsection{The shifted-Laplacian response \texorpdfstring{$\Cph$}{C\_phi}}\label{ssec:cphi}

The substrate's linear response to an external source $f \in \mathbb{R}^{120}$
is the discrete Green's function of the shifted Laplacian:
\begin{equation}\label{eq:cphi}
\Cph \;=\; \Lop + \Ph^{-2} I,
\qquad
\psi \;=\; \Cph^{-1} f.
\end{equation}
The shift $\Ph^{-2} \approx 0.382$ is a design-level stability
clamp: it makes $\Cph$ strictly positive definite (smallest eigenvalue
$\Ph^{-2}$, since the trivial Laplacian eigenvalue is $0$), so that the
inverse is bounded with operator norm $\Ph^{2}\approx 2.618$. This is
\emph{not} a derived theorem; it is a stability choice. The companion
kernel document~\citep{SmartAriaClosureKernel2026} discusses the same
$\Cph$ as the basis for an independent passive-regime witness in
flavour physics~\citep{SmartBAnomaly2026}, where the same operator
form (without retuning the shift) describes the
$B\to K^{*}\mu^{+}\mu^{-}$ angular anomaly across five public datasets.
This paper imports $\Cph$ from that line; we do not re-derive it.

The response $\psi = \Cph^{-1} f$ is smooth in $f$. By itself it does
not produce critical-state cascade statistics; the recurrent layer
in~\S\ref{sec:chain} adds a graph-pinned nonlinearity
$\mathrm{bounded\_topk}(\psi, k\!=\!12)$ to obtain self-organised-critical
event distributions. The choice $k\!=\!12$ is the average degree of
$G_{V_{600}}$ (\S\ref{ssec:graph}), pinned by the substrate, not
fitted to any dataset.

\paragraph{Disclosure (substrate-witness scope).}
The $\Ph^{-2}$ floor is a stability shift, not a derived parameter.
The bounded-top-$K$ nonlinearity uses $k\!=\!12$ pinned to the graph's
average degree, not a fitted threshold. No other shape parameter
enters. The condition-dependent self-injection coupling
$\eta\in\{0, 0.05, 0.20\}$ is the only architectural parameter that
varies between conditions in~\S\ref{sec:chain}; it is reported
explicitly per condition.

\subsection{Shell decomposition}\label{ssec:shells}

Under the H$_3$ subgroup, the $120$ vertices partition into $9$
spherical shells indexed by Euclidean inner product with a chosen pole:
\[
\bigl(|S_0|,\ldots,|S_8|\bigr) \;=\; (1, 12, 20, 12, 30, 12, 20, 12, 1).
\]
The middle shell $S_{4}$ has $30$ vertices on the equatorial plane
(the icosidodecahedral ring). When projecting onto a continuum kernel
in companion preprints~\citep{SmartAriaClosureKernel2026}, the
shell-mean projection of the equatorial-source response,
$\kappa(x) = \mathrm{shell\text{-}mean}[\Cph^{-1} f_{\mathrm{equat}}](x)$,
collapses on the continuum exponential $\frac{\Ph}{2}\,e^{-|x|/\Ph}$.
This paper does not use that continuum projection; we work with the
discrete operator throughout.

\subsection{Cascade descent (sketch)}\label{ssec:cascade}

The recurrent layer in~\S\ref{sec:chain} references a cascade
decomposition $E_{8}\to H_{4}\to H_{3}\to D_{4}\to \mathrm{Cl}(1,3)
\to S^{7}\to 0$, implemented in
\texttt{kernel/cascade\_descent.py:descend\_operator\_e8\_to\_h4\_clean}.
An arbitrary operator on the $E_{8}$ root system descends to the
4-D H$_4$ subspace through a sequence of orthogonal projections, each
preserving the Frobenius norm to within $10^{-15}$. The
$\sigma$-orbit projector basis from
\texttt{kernel/sigma\_orbit\_basis.py} gives the K-class block
decomposition at machine precision.

The descent provides numerical stability for the cascade ablations:
when one ablates a specific K-class contribution (e.g.\ $K_{7}$), the
remaining operator structure is exactly preserved. We do not claim
the cascade itself is forced by physics on a pre-substrate level; the
cascade is a decomposition of operators on $\Rsixhundred$, and the
choice of $\Rsixhundred$ as the active substrate is post-hoc justified
by the empirical correspondences in~\S\ref{sec:results}.

\subsection{What the substrate is fixed-by, and what it is not}

\begin{itemize}\itemsep=2pt
\item Fixed by group theory once $\Rsixhundred$ is chosen: $|V|=120$,
  uniform degree $12$, Laplacian spectrum, $9$-shell partition, K-class
  irrep block structure, average degree $k\!=\!12$, $\sigma$-twin pairing.
\item Fixed by stability choice: the $\Ph^{-2}$ shift in $\Cph$. This
  is not a derivation; it is a design-level clamp that bounds the
  response inverse.
\item Not fixed by this paper: the choice of $\Rsixhundred$ over the
  $24$-cell or $120$-cell. The choice is post-hoc motivated by the H$_4$
  cascade structure and the empirical correspondences. A formal
  ablation against alternative regular 4-polytopes is an open build
  (\S\ref{sec:limitations}).
\end{itemize}
